\documentclass[10pt,a4paper]{amsart}
\usepackage[margin=19mm]{geometry}
\usepackage{amsmath,amssymb,mathtools}
\usepackage[scr=rsfs,cal=euler]{mathalfa}
\usepackage{xfrac}
\usepackage[unicode,hidelinks,bookmarks=false]{hyperref}
\newcommand{\ie}{i.e.\ }
\newcommand{\cf}{cf.\ }
\newcommand{\resp}{respectively\ }
\newcommand{\Subsection}{\subsection}
\providecommand{\nobreakdash}{\nobreak\hbox{-}}
\setlength{\emergencystretch}{2em}
\multlinegap0pt
\usepackage{bm}
\usepackage{bbm}
\usepackage{dsfont}
\usepackage{xspace}
\usepackage{cancel}
\usepackage{mathtools}
\usepackage{amsthm}
\makeatletter
\@ifundefined{theorem}{\newtheorem{theorem}{Theorem}}{}
\@ifundefined{lemma}{\newtheorem{lemma}[theorem]{Lemma}}{}
\makeatother
\renewcommand{\geq}{\geqslant}
\renewcommand{\leq}{\leqslant}
\title[$2$-designs admitting a flag-transitive automorphism group w]{$2$-designs admitting a flag-transitive automorphism group with socle $PSL(2,q)$}
\author{Hongxue Liang, Mario Galici, Zhihui Liu, Filip Martinovi\'c, Alessandro Montinaro, Eleonora Romano}
\date{Source: arXiv:2607.05067v1 (CC BY 4.0)}
\begin{document}
\maketitle

\noindent\emph{Source and licence.} $2$-designs admitting a flag-transitive automorphism group with socle $PSL(2,q)$. Version arXiv:2607.05067v1 (CC BY 4.0), distributed under the Creative Commons Attribution 4.0 International licence, as recorded in the arXiv licence metadata for this exact version and shown on the abstract page https://arxiv.org/abs/2607.05067v1. Full rights notice: licenses/arxiv-2607.05067-CC-BY-4.0.txt.\par\smallskip

\noindent\textit{Editorial scope.} Counted unit: the Lemma whose source label is \texttt{OndaDobro} in the article (source0.tex lines 2861--2869, source label \texttt{OndaDobro}), delivered together with the article's own complete original proof of it (source0.tex lines 2871--2934, one \texttt{proof} environment of 64 lines). No mathematics is written, repaired or re-derived: statement and proof are the article's own text line for line. Required context retained uncounted: lem:basic-params (lines 554--562); SigmaMax1 (lines 2804--2819). Every definition, display and label of those lines is retained unchanged. Omitted from this package and not redistributed: 1--553, 563--2803, 2820--2860, 2870--2870, 2935--4884. Nothing inside a retained excerpt is omitted or altered: every retained body is delivered exactly as the article's lines are written, so no omission receipt of any kind accompanies this package. Citations: the retained excerpts carry no citation command at all, so no bibliography entry of the article is transcribed and no reference is redistributed. Reference adaptation: none was needed and none was made; every reference inside the delivered unit resolves inside it, so no unresolved reference or citation is produced. Printed slips kept literal: no printed word, symbol, hypothesis or number of the counted statement or of its proof was altered. Layout and class: the article's own class is \texttt{amsart} with packages including rotating, eurosym, geometry, tabularx, graphicx, longtable, makecell, booktabs, threeparttable, array, amsmath, amssymb, amsfonts, amsthm; the wrapper is the lane's pinned \texttt{amsart} editorial wrapper at 10pt on A4, which restates the theorem-like environments the retained text needs and adds the article's own macro definitions that the retained text needs (\texttt{\textbackslash geq}, \texttt{\textbackslash leq}), with extra wrapper packages (bm, bbm, dsfont, xspace, cancel, mathtools). The delivered numbering is the unit's own. Assets: no retained span contains a figure environment, a graphics-inclusion command, a PDF-page inclusion, a table or a TikZ picture; no image file is copied into this package, the package declares no asset dependency and no image file is redistributed with it. Licence: version arXiv:2607.05067v1 of this article is distributed under the Creative Commons Attribution 4.0 International licence, as recorded in the arXiv licence metadata for this exact version and shown on the abstract page https://arxiv.org/abs/2607.05067v1. A full rights notice is delivered at licenses/arxiv-2607.05067-CC-BY-4.0.txt. Characters: every retained excerpt is pure ASCII, so the delivered engine, XeLaTeX with its default Latin Modern text font, reports no missing character.
\begin{lemma} \label{lem:basic-params}
Let $\mathcal{D}$ be a $2$-$(v,k,\lambda)$ design admitting a flag-transitive
automorphism group $G$. Let $\alpha \in \mathcal{P}$, then
\begin{enumerate}
\item[\rm(i)]   $r(k-1)=\lambda(v-1)$, $vr=bk$ and $b \geq v$;
\item[\rm(ii)]  $r\mid \lambda \gcd\left(\left\vert G_\alpha\right\vert, v-1\right)$, and $r^{2}>\lambda v$;
\item[\rm(iii)]  $r\mid \lambda d$ for every nontrivial subdegree $d$ of $G$.
\end{enumerate}
\end{lemma}

\begin{lemma}
\label{SigmaMax1}If $\Sigma $ is any maximal $G$-invariant partition, then

\begin{enumerate}
\item $v_{0}=p^{a}\theta +1$ with $\theta =\left( k_{0}-1,v_{0}-1\right) $
and $1\leq a\leq f$.

\item $r_{0}=\frac{p^{a}\lambda _{0}}{\left( k_{0}-1\right) /\theta }$;

\item $k_{0}\mid v_{0}\left( p^{f-a},\frac{p^{a}\lambda _{0}}{\left(
k_{0}-1\right) /\theta }\right) $;

\item $k_{1}=\frac{k_{0}-1}{\theta }\cdot \frac{p^{f-a}v_{0}}{k_{0}}+1$ and $%
k_{1}>\frac{p^{f}(k_{0}-1)}{k_{0}}+1$.
\end{enumerate}
\end{lemma}

\begin{lemma}
\label{OndaDobro}If $\Sigma $ is any maximal $G$-invariant partition, and $\Delta \in \Sigma$, then the following hold:
\begin{enumerate}
    \item $G_{\Delta }^{\Delta }=Z_{v_{0}}:Z_{d}$, with $v_{0}=\frac{p^{f}-1}{%
\left\vert X_{(\Delta )}\right\vert _{p^{\prime }}(2,p^{f}-1)}$ and $d\mid f$;
  \item  $v_{0}=p^{a}\theta +1$ with $\theta =\left( k_{0}-1,v_{0}-1\right) $
and $1\leq a <f$.
\end{enumerate}
\end{lemma}

\begin{proof}
The group $X$ acts $2$-transitively on the $v_{1}=q+1$ points of $\mathcal{D}_{1}$, hence $X_{\Delta }\cong Z^{f}_{p}:Z_{\frac{p^{f}-1}{(2,p^{f}-1)}}$. Furthermore, $X_{\Delta }\trianglelefteq G_{\Delta }\leq A\Gamma L(1,p^{f})$. It follows from Lemma \ref{SigmaMax1}(1) that $(v_{0},p)=1$, hence $(Z_{p})^{f}\leq X_{x}$ for some $x\in \Delta $. Then $(Z_{p})^{f}\leq X_{(\Delta )}$ since $(Z_{p})^{f}\trianglelefteq G_{\Delta }$ and $G_{\Delta }$ acts
transitively on $\Delta $.

If $X_{(\Delta )}=Z_{p}^{f}:Z_{\frac{p^{f}-1}{(2,p^{f}-1)}}$, then $%
G_{\Delta }^{\Delta }$ is isomorphic to a subgroup of $Z_{f}$. Then $%
G_{\Delta }^{\Delta }$ is abelian, and hence $G_{\Delta }^{\Delta }$ is
regular since $G_{\Delta }^{\Delta }$ acts transitively on $\Delta $,
however this is impossible since $r_{0}\mid \left\vert G_{x}^{\Delta
}\right\vert $ with $r_{0}>1$. Thus $Z_{p}^{f}\leq X_{(\Delta )}<
(Z_{p})^{f}:Z_{\frac{p^{f}-1}{(2,p^{f}-1)}}$, and hence $X_{\Delta }^{\Delta
}\cong Z_{c}$, with $c>1$ and $c\left\vert X_{(\Delta )}\right\vert _{p^{\prime }}=%
\frac{p^{f}-1}{(2,p^{f}-1)}$, acts point-semiregularly on $%
\mathcal{D}_{0}$. Since $G_{\Delta }^{\Delta }$ acts transitively on $\Delta 
$ and $1\neq X_{\Delta }^{\Delta }\trianglelefteq G_{\Delta }^{\Delta }$, it
follows that $c\mid v_{0}$. Moreover, $G_{\Delta }^{\Delta }/X_{\Delta
}^{\Delta }\cong Z_{d}$, with $d\mid f$, permutes transitively the $\frac{v_{0}%
}{c}$ orbits on $\Delta $ under $X_{\Delta }^{\Delta }$. Since $G_{\Delta
}^{\Delta }$ acts flag-transitively on $\mathcal{D}_{0}$, it follows that $%
r_{0}\mid \left\vert G_{x}^{\Delta }\right\vert $. On the other hand, $%
G_{x}^{\Delta }\cap X_{\Delta }^{\Delta }=1$ since $X_{\Delta }^{\Delta }$
acts semiregularly on $\Delta $. Therefore, $G_{x}^{\Delta }$ is isomorphic
to a subgroup of $Z_{d}$, and $r_{0}\mid d$.

Note that $X_{\Delta
}^{\Delta }:G_{x}^{\Delta }\trianglelefteq G_{\Delta }^{\Delta }$ since $G_{\Delta }^{\Delta }/X_{\Delta }^{\Delta }\cong Z_{d}$. Then $\Delta $ is a
union of $(X_{\Delta }^{\Delta }:G_{x}^{\Delta })$-orbits of equal lengths.
On the other hand, $x^{X_{\Delta }^{\Delta }}$ is both a $X_{\Delta
}^{\Delta }$-orbit and a $(X_{\Delta }^{\Delta }:G_{x}^{\Delta })$-orbit.
Thus, each $X_{\Delta }^{\Delta }$-orbit on $\Delta $ has length $c$ and is a $(X_{\Delta
}^{\Delta }:G_{x}^{\Delta })$-orbit. In particular, the $x^{X_{\Delta
}^{\Delta }}\setminus \{x\}$ is a union of $G_{x}^{\Delta }$-orbits. If $C$
is any block of $\mathcal{D}_{0}$ containing $x$, then 
\begin{equation*}
\frac{\left( v_{0}-1\right) \lambda _{0}}{k_{0}-1}\left\vert C\cap
(x^{X_{\Delta }^{\Delta }}\setminus \{x\})\right\vert =\left\vert
x^{X_{\Delta }^{\Delta }}\setminus \{x\}\right\vert \lambda _{0}\text{,}
\end{equation*}
by Lemma \ref{lem:basic-params}(iii) with $\mathcal{D}_{0}$ and $G_{\Delta }^{\Delta }$ in the role of $\mathcal{D}$ and $G$, respectively. So, we obtain 
%$$\left( v_{0}-1\right) \left\vert C\cap (X_{\Delta }^{\Delta }\setminus
%\{x\})\right\vert =(c-1)(k_{0}-1)\text{,}$$
%{\color{blue}
$$\left( v_{0}-1\right) \left\vert C\cap (x^{X_{\Delta }^{\Delta }}\setminus
\{x\})\right\vert =(c-1)(k_{0}-1)\text{,}$$
which implies $v_{0}-1 \mid (c-1)(k_{0}-1)$.

If $G_{\Delta }^{\Delta }\neq X_{\Delta }^{\Delta
}:G_{x}^{\Delta }$, then there is $y\in \Delta \setminus x^{X_{\Delta
}^{\Delta }}$, and hence $y^{X_{\Delta }^{\Delta }}$ is a $(X_{\Delta
}^{\Delta }:G_{x}^{\Delta })$-orbit. Then $y^{X_{\Delta }^{\Delta }}$ is a
union $G_{x}^{\Delta }$-orbits and so%
\begin{equation*}
\frac{\left( v_{0}-1\right) \lambda _{0}}{k_{0}-1}\left\vert C\cap
y^{X_{\Delta }^{\Delta }}\right\vert =\left\vert y^{X_{\Delta }^{\Delta
}}\right\vert \lambda _{0}\text{,}
\end{equation*}%
by Lemma \ref{lem:basic-params}(iii), which implies $v_{0}-1 \mid c(k_{0}-1)$. Therefore $v_{0}-1\mid k_{0}-1$, since we have seen that $v_{0}-1 \mid (c-1)(k_{0}-1)$, which is a contradiction. Thus $%
G_{\Delta }^{\Delta }=X_{\Delta }^{\Delta }:G_{x}^{\Delta }$ and $X_{\Delta
}^{\Delta }$ acts regularly on $\Delta $. In particular $c=v_{0}$ and $%
\left\vert G_{x}^{\Delta }\right\vert \mid f$. Hence, $G_{\Delta }^{\Delta
}=Z_{v_{0}}:Z_{d}$ \ with $d\mid f$, which is the assertion (1).

Thus $v_{0}=\left\vert G_{\Delta }:G_{x}\right\vert=\left\vert X_{\Delta }:X_{(\Delta )}\right\vert $, and hence $v_{0}=\frac{p^{f}-1}{\left\vert X_{(\Delta)}\right\vert _{p^{\prime }}(2,p^{f}-1)}$. Then $a<f$ since $v_{0}=p^{a}\theta+1$ with $\theta =\left( k_{0}-1,v_{0}-1\right)$ and $1\leq a \leq f$ by Lemma \ref{SigmaMax1}(1), and we obtain the asserton (2). 
\end{proof}

\end{document}
